\noindent These tables list the reduced basis elements and involution signs by sigma number.
\begin{tableblock}
\captionsetup{type=table,font=normalsize,labelfont=bf,skip=4pt}
\captionof{table}{Sigma basis elements by sigma number.}\label{tab:sigma-number-components}
\renewcommand{\arraystretch}{1.12}
\setlength{\tabcolsep}{4pt}
\normalsize
\begin{tabularx}{\columnwidth}{@{}l >{\raggedright\arraybackslash}X@{}}
\toprule
sigma number & representative basis elements \\
\midrule
\(0\) (scalar) & \(1\) \\
\(1\) (vector) & \(\{\sigma_1,\ \sigma_2,\ \sigma_3\}\) \\
\(2\) (bivector) & \(\{\sigma_2\sigma_3,\ \sigma_3\sigma_1,\ \sigma_1\sigma_2\}\) \\
\(3\) (pseudoscalar) & \(\sigma_1\sigma_2\sigma_3\) \\
\bottomrule
\end{tabularx}
\end{tableblock}

\begin{tableblock}
\captionsetup{type=table,font=normalsize,labelfont=bf,skip=4pt}
\captionof{table}{Involution signs by sigma number.}\label{tab:sigma-number-signs}
\renewcommand{\arraystretch}{1.12}
\setlength{\tabcolsep}{4pt}
\normalsize
\begin{tabularx}{\columnwidth}{@{}l >{\raggedright\arraybackslash}X >{\raggedright\arraybackslash}X >{\raggedright\arraybackslash}X@{}}
\toprule
sigma number & \((\cdot)^*\) & \((\cdot)^{\mathrm R}\) & \((\cdot)^{*\mathrm R}\) \\
\midrule
\(0\) (scalar) & \(+\) & \(+\) & \(+\) \\
\(1\) (vector) & \(-\) & \(+\) & \(-\) \\
\(2\) (bivector) & \(+\) & \(-\) & \(-\) \\
\(3\) (pseudoscalar) & \(-\) & \(-\) & \(+\) \\
\bottomrule
\end{tabularx}
\end{tableblock}
